% Same-content excerpts from the user-supplied Adobe/Word PDF, without annotations.
\documentclass[degree=master,degreetype=profession]{tongjithesis}
\tongjisetup{ctitle={学位论文中文题目},etitle={Thesis Title},cauthor={},eauthor={},studentnumber={},cdepartment={},edepartment={},csubjectcategory={},esubjectcategory={},cmajorfirst={},emajorfirst={},cmajorsecond={},emajorsecond={},cresearchfield={},eresearchfield={},csupervisor={},esupervisor={},cassosupervisor={\mbox{}},eassosupervisor={\mbox{}},cjointinstitution={\mbox{}},ejointinstitution={\mbox{}},cdate={二〇XX年XX月},edate={Month, Year},edegree={Master of Engineering}}
\begin{cabstract}
在实际工程结构的服役过程中，由于非线性与随机性的耦合作用，工程结构特别是混凝土结构的非线性反应具有不可精确预测的性质。因此，从概率密度演化的角度考察工程结构的非线性性状是准确把握结构非线性性能的必由之路。本文基于随机结构反应概率密度演化的思想对于随机结构分析理论进行了深入的探讨，初步建立了随机结构反应概率密度演化的基本图景。
\end{cabstract}
\ckeywords{随机结构,马尔可夫过程,非线性构形状态,差分方法}
\begin{eabstract}This page is not part of the visual comparison.\end{eabstract}
\ekeywords{example}
\begin{document}
\graphicspath{{figures/}}
\frontmatter\makecover
\mainmatter
\chapter{引言}\section{概述}
随着现代科学技术的发展和人们生活质量的提高，对于工程结构的性能提出了越来越高的要求。例如：现代精密仪器、大型设备往往对于振动与位移有严格的限制；生命线工程结构，要求在大震和大灾作用下依然保有必要的功能，以为灾后救援与重建提供保障。20世纪中叶以来，尽管社会发展水平有了巨大的提高，然而由于灾害性作用而造成的损失却反而越来越大，这给结构工程学科带来了一系列新的挑战性课题。正是在这样的背景下，基于性能的设计思想开始浮出水面，并在近十年来引起了学者们强烈的兴趣。
\makeatletter
\def\tongji@cheadingtitlenew{{\fontspec{SimSun.ttf}[Path=fonts/]X\allowbreak X\allowbreak X\allowbreak X\allowbreak X\allowbreak X\allowbreak X\allowbreak X\allowbreak X\allowbreak X\allowbreak X\allowbreak X\allowbreak X\allowbreak X\allowbreak X\allowbreak X\allowbreak X\allowbreak X\allowbreak X\allowbreak X\allowbreak X\allowbreak X\allowbreak X\allowbreak X\allowbreak X\allowbreak X\allowbreak X\allowbreak X\allowbreak X\allowbreak X\allowbreak X\allowbreak X\allowbreak X\allowbreak X\allowbreak X\allowbreak X\allowbreak X\allowbreak X\allowbreak X\allowbreak X\allowbreak X\allowbreak X\allowbreak X\allowbreak X\allowbreak X\allowbreak X\allowbreak X\allowbreak X\allowbreak X\allowbreak X\allowbreak X\allowbreak X\allowbreak X\allowbreak X\allowbreak X\allowbreak X\allowbreak X\allowbreak X\allowbreak X\allowbreak X\allowbreak X\allowbreak X\allowbreak X\allowbreak X\allowbreak X\allowbreak X\allowbreak X\allowbreak X\allowbreak X\allowbreak X\allowbreak X\allowbreak X\allowbreak X\allowbreak X\allowbreak X\allowbreak X\allowbreak X\allowbreak X\allowbreak X\allowbreak X\allowbreak X\allowbreak }}
\makeatother
\begin{statements}\input{data/statements}\end{statements}
\end{document}
